% Rigour check of the unified-model work of 15 October 2026 (route 2 of P3_OUTLINE.ja.md).
% Build: pdflatex rigor_check.tex (twice).  Companion of unified_model_derivation.tex.
\documentclass[11pt,a4paper]{article}
\usepackage[margin=2.1cm]{geometry}
\usepackage{amsmath,amssymb}
\usepackage{mathptmx}
\usepackage{booktabs,tabularx,array}
\usepackage{enumitem}
\usepackage[hyphens]{url}
\urlstyle{tt}
\setlist{nosep,leftmargin=*}
\newcolumntype{Y}{>{\raggedright\arraybackslash}X}
\begin{document}
\begin{center}
{\large\bfseries Rigour check: the two-level biofilm model from one variational principle}\\[3pt]
Keisuke Nishioka, working note, 15 October 2026 (not part of the thesis)
\end{center}

\noindent This note lists, for the model of \path{unified_model_derivation.tex}, what is
derived, what is checked by machine, what is checked by computation, what is assumed, and
what is still open. Every entry names the file that produces it. All files are in
\texttt{JAXFEM/} of the repository unless stated otherwise.

\section{Derivation}

\small\begin{tabularx}{\linewidth}{@{}Yp{2.9cm}Y@{}}
\toprule
Statement & Status & How it is established\\
\midrule
The discrete equations solved are the Euler--Lagrange equations of the discrete functional
$E_h+\Delta t\,D_h+C_h$ & established & The residual is the automatic derivative (JAX) of the
functional; no equation is written by hand (\path{unified_two_level_1d.py},
\texttt{functional}, \texttt{residual}).\\
The point model of Klempt et al.\ 2026 is the stationarity condition of the energy
$\Psi_{\rm pm}$ and dissipation $D_{\rm pm}$ used here & established & One node, no coupling:
the composition rows equal the repository's point-model residual
(\path{core_hamilton_1d.residual}) to round-off, every row
(\path{unified_two_level_variational_check.py}).\\
The hand-written residual of the first version agrees with the automatic one & established,
with one known difference & Composition rows to round-off; the phase row differs on purpose
(the hand version approximated one term and let the barrier terms act on the phase).\\
Scheme (c) of the element is the limit $\lambda\to0$, $\beta_s=0$ of the model & established
by inspection & The composition equations lose the coupling terms and the expansion equation
becomes $\dot\alpha=k_\alpha\phi\sum_i s_i\psi_i$ (derivation note, Section~4).\\
$\dfrac{\mathrm d}{\mathrm dt}\int E=-2\int D\le0$ without external power & proved for the
model with the barrier terms as stored energy & Derivation note, Section~5. For the version
with the barrier terms as a constraint the inequality holds in the limit $K\to0$; a discrete
check is open.\\
The front motion follows from the growth kinematics & established in 1D, 2D and 3D by
computation & Below.\\
\bottomrule
\end{tabularx}\normalsize

\section{Checks by computation}

\small\begin{tabularx}{\linewidth}{@{}YYp{4.2cm}@{}}
\toprule
Check & Result & File\\
\midrule
Uniform field, no coupling: every node equals the point model solved to convergence &
$8.9\times10^{-11}$ (two-level), $1.6\times10^{-13}$ (one-level) & \path{unified_two_level_1d.py},
\path{unified_1d_monolithic.py}\\
Energy decrease of the one-level model, every step & largest change $-5.6\times10^{-3}$,
no increase except at the antibiotic switch (external input) & \path{unified_1d_monolithic.py}\\
Energy decrease of the fuel model with closed boundary & largest change
$-9.8\times10^{-3}$, no increase & \path{unified_two_level_fuel_1d.py}\\
Newton convergence & residual $<10^{-8}$ every step, at most 31 iterations, adaptive
halving of the step when needed & \path{unified_two_level_1d.py}\\
Time step, two-level 1D & 60 against 120 steps: $\alpha-1$ changes by $0.7\,\%$ &
\path{unified_two_level_1d.py}\\
Mesh, 3D lumped & $16^3$, $24^3$, $32^3$: edge displacement within $3\,\%$ &
\path{fuel_front_3d.py}\\
Mesh, 3D five species & $16^3$ against $24^3$: within $8\,\%$ (constant $c^*$),
$9\,\%$ (local nutrient) & \path{fuel_front_3d_species.py}\\
Point model inside the 3D runs solved to convergence & residual $\le4.5\times10^{-13}$
every step & \path{fuel_front_3d_species.py}\\
\bottomrule
\end{tabularx}\normalsize

\section{Assumptions, each with its reason}

\small\begin{tabularx}{\linewidth}{@{}YYp{2.4cm}@{}}
\toprule
Assumption & Reason & Tested\\
\midrule
Two levels (phase $\phi$ and composition $s_i,s_0,\psi_i$) & One level fills the outer liquid
from the barrier floor (\path{unified_1d_monolithic.py}); the two voids are different
objects & yes\\
Weight $h(\phi)=\phi^2$ of the composition energy & With $h=\phi$ the composition force does
not vanish in the outer liquid and it fills again; $\phi^2$ is the weight of the mechanical
term of Klempt et al.\ 2024 & yes\\
The barrier terms act as a constraint, not as stored energy & Their value (517 at the
calibration's start against 8 at $\psi=0.5$) would make dying cells grow the biofilm & yes\\
$\phi\le1$ by a one-sided barrier $K_\phi/(1-\phi)^2$ & The phase equation has no upper
bound (Klempt et al.\ 2024 neither); $\phi$ reached 2.7 at large $k_\alpha$ & yes\\
Growth as a nutrient reaction (Onsager), the living biomass as catalyst & A growth energy on
the living biomass keeps $\psi$ up at large growth and defeats the antibiotic; the catalyst form
has no force on $\psi$ and makes the total energy decrease & yes\\
\bottomrule
\end{tabularx}

\medskip\noindent\small\begin{tabularx}{\linewidth}{@{}YYp{2.4cm}@{}}
\toprule
Assumption & Reason & Tested\\
\midrule
Ideal-solute free energy $c(\ln c-1)$ of the nutrient, linear kinetics
$\dot\xi=L(Y\ln c-\mu_b)$ & Simplest forms consistent with the principle; $\mu_b=Y\ln c_{\rm stop}$
gives a stop concentration & sensitivity only\\
$\lambda$ (energy scale of the composition), $\beta_s$, $L_0$, $Y$, $c_{\rm stop}$, $D$ &
Not from a paper; $\lambda$ is the one that matters for the antibiotic's effect on the phase,
$\beta_s$ does not matter for a seed & $\lambda$, $\beta_s$: yes; others: no\\
Time link of the composition ($s=0.2375$ per $T^*$) & The calibration's link; in the model it
is the ratio of dissipation coefficients and can be fixed by the measured volume & no\\
Local nutrient in the point model: interaction $c^*c(\mathbf x)/c_{\rm ref}$, $c_{\rm ref}=1$ &
The element's local-nutrient option; the relation of $c^*$ to the field's $c$ should come from
the nutrient level of the calibration & no\\
Front term of Klempt et al.\ 2024 left out & Not derivable from a potential (their Remark~1);
the edge moves towards the nutrient by the growth kinematics alone & yes (1D, 2D, 3D)\\
Mechanics in the 2D/3D checks: small strain, $\nu=0.3$, linear elasticity, nutrient in the
reference configuration & Checks of the kinematics only; the full model uses Klempt et al.\
2024's mechanics ($\nu=0.49$) & no\\
Lumped composition in \path{fuel_front_2d.py}, \path{fuel_front_3d.py} & Kinematics check;
replaced by the five-species point model in \path{fuel_front_3d_species.py} & yes\\
\bottomrule
\end{tabularx}\normalsize

\section{Open}

\begin{enumerate}
\item A discrete energy check of the two-level model with the barrier terms as a constraint.
\item The limit $\phi\to0$ of the composition equations (regularised by $\varepsilon$).
\item The composition acting on the phase ($\lambda>0$) and the composition diffusion
  ($\beta_s$) in 3D; the 3D runs so far use the decoupled limit.
\item Finite strain and $\nu=0.49$ in the 2D/3D checks; nutrient diffusion in the current
  configuration.
\item Existence and uniqueness of solutions (not attempted).
\item A reader other than me: the derivation has not been checked by Dr.\ Soleimani or
  anyone else.
\item Overlap with F.\ Klempt's unpublished two-species work: only published sources were
  used; to be discussed before any publication.
\end{enumerate}

\section{What the checks do not cover}

The checks establish that the equations are the ones the energy and dissipation imply, that
the solver solves them, and that the results do not depend on the mesh or the time step at
the level stated. They do not establish that the model is right: the choices in Section~3 are
modelling assumptions, and the only data used so far are the calibrated interactions of the
point model. The measured volumes (Heine et al.) and an antibiotic experiment would test the
growth law and $\lambda$.
\end{document}
